\chapter{总线、地址译码与片上互连}

\begin{keyidea}
互连的核心不是一束导线，而是请求、目标选择、背压、响应和错误的事务契约。本章从地址地图进入总线与片上互连。
\end{keyidea}

\section{一、地址译码定义整机地图}

CPU 输出的是地址数值，整机必须规定这个数值由哪个目标响应。这个规定叫地址映射。一个简单教学板可以把低地址映射到 ROM，中间映射到 RAM，高地址小窗口映射到 GPIO、定时器和串口等设备寄存器。地址译码器把高位地址和访问类型变成片选信号。

\begin{lstlisting}
example 16-bit address map:
  0x0000-0x7FFF  boot ROM
  0x8000-0xEFFF  SRAM
  0xF000         GPIO_OUT
  0xF004         GPIO_IN
  0xF010         TIMER_COUNT
  0xF014         TIMER_CTRL
\end{lstlisting}

\begin{pdfdiagram}
  \node[hw state, minimum width=2.0cm] (addr) at (0,0) {地址 A[15:0]};
  \node[hw compute, minimum width=2.15cm] (dec) at (3.0,0) {高位译码};
  \node[hw state, minimum width=1.6cm] (rom) at (6.0,1.05) {ROM CS};
  \node[hw state, minimum width=1.6cm] (ram) at (6.0,0) {RAM CS};
  \node[hw control, minimum width=1.6cm] (io) at (6.0,-1.05) {IO CS};
  \draw[hw bus] (addr) -- (dec);
  \draw[hw bus] (dec.east) |- (rom.west) node[pos=0.3,above,hw label]{A15=0};
  \draw[hw bus] (dec) -- (ram) node[pos=0.55,above,hw label]{1xxx（非 1111）};
  \draw[hw bus] (dec.east) |- (io.west) node[pos=0.3,below,hw label]{1111};
  \node[hw label, text width=7.8cm] at (3.4,-1.95) {设计目标：任意一个地址最多选中一个目标；未映射区域也要有错误、固定值或超时规则。};
\end{pdfdiagram}
\diagramnote{地址译码不是软件表格，而是片选电路。片选重叠会造成总线争用或误写设备。}

\begin{longtable}{L{0.18\linewidth}L{0.38\linewidth}L{0.34\linewidth}}
\toprule
映射对象 & 响应规则 & 典型错误 \\
\midrule
ROM & 复位入口和只读代码地址命中 & 复位后 CPU 取不到第一条指令 \\\tablerowrule
RAM & 普通数据读写地址命中 & 写使能误打到 ROM 或设备窗口 \\\tablerowrule
GPIO & 设备寄存器地址命中并产生副作用 & 地址译码过宽，普通内存写误触发外设 \\\tablerowrule
未映射区 & 返回错误、超时或固定空值 & 访问悬空，CPU 永久等待或读到随机总线值 \\
\bottomrule
\end{longtable}

\topic{片选表达式就是硬件约定}
若地址空间为 64 KiB，可以用高位地址产生粗粒度片选。例如 \code{A15=0} 选 ROM，\code{A15=1} 且 \code{A14,A13,A12} 不全为 1 选 RAM（\code{0x8000-0xEFFF}，28 KiB），\code{A15=1,A14=1,A13=1,A12=1} 选 I/O 窗口。无论用门电路、译码器、CPLD 还是 FPGA 实现，都必须证明任意合法周期最多只有一个目标片选有效。

\begin{lstlisting}
rom_cs = !A15
ram_cs = A15 && !(A14 && A13 && A12)
io_cs  = A15 && A14 && A13 && A12
\end{lstlisting}

这三个表达式看起来像逻辑练习，实际上就是整机安全边界。若 ROM 和 RAM 可同时片选，读周期可能出现两个芯片同时驱动数据线；若 RAM 和 GPIO 片选重叠，普通 STORE 可能既改内存又改外设输出；若未映射地址没有超时或错误返回，CPU 可能等待一个永远不会响应的目标。

\topic{时序预算决定是否需要等待状态}
地址译码需要时间，存储器访问需要时间，数据到达 CPU 输入端后还要满足建立时间。若这些时间总和超过时钟周期，就必须降低频率、换更快器件，或插入等待状态。等待状态不是软件延时，而是硬件事务延长：地址、片选、读写方向和数据驱动关系必须在等待期间保持可解释。

\begin{lstlisting}
read_budget =
  cpu_address_output_delay
  + address_decode_delay
  + memory_access_time
  + input_setup_time
  + board_margin
\end{lstlisting}

若 \code{read_budget} 大于一个周期，控制器就不能在下一拍采样数据。慢 ROM、慢外设和长板级走线都可能需要 READY/WAIT 机制。正式规格还应说明最大等待长度和超时行为，避免坏设备把整机永久挂住。

\begin{longtable}{L{0.2\linewidth}L{0.34\linewidth}L{0.36\linewidth}}
\toprule
处理方式 & 适合场景 & 代价 \\
\midrule
降低时钟 & 全系统都慢，教学板容易观察 & 所有操作都变慢 \\\tablerowrule
更快器件 & SRAM 或译码器成为瓶颈 & 成本、功耗或电平兼容性变化 \\\tablerowrule
插入等待 & 只有部分目标较慢 & 控制器必须保持事务状态 \\\tablerowrule
超时返回错误 & 目标可能不存在或失效 & 需要错误状态和恢复入口 \\
\bottomrule
\end{longtable}

\topic{案例：用 74HC138 译出整机片选}
本节开头给出了 16-bit 地图：\code{0x0000-0x7FFF} 是 boot ROM（32 KiB，A15=0），\code{0x8000-0xEFFF} 是 SRAM（28 KiB），\code{0xF000} 起是 I/O 区（4 KiB）。现在给出这张地图的引脚级实现。最经典的器件是 3-8 译码器 74HC138：它有三个地址输入 C、B、A 和八个低有效输出 Y0--Y7，任一时刻只有一个输出为低。把 A[15:13] 依次接上 C、B、A，每个输出就恰好管辖一个 8 KiB 窗口，窗内偏移由 A[12:0] 决定；使能端 G1 接高电平、/G2A 与 /G2B 接低电平，译码器才工作——这两个使能端也常改接总线读写选通，让片选只在存储周期内出现。输出低有效带来一个方便的性质：把若干相邻窗口合并成一个片选只需一个与门，任何一个被选中的输出变低，与门输出即为低。于是 Y0--Y3 经四输入与门合成 ROM\_CS，正好覆盖 A15=0 的整个低半空间；Y4--Y6 经三输入与门合成 RAM\_CS；Y7 则送入二级译码，由 A12 把 \code{0xE000-0xFFFF} 再分成 SRAM 尾窗和 I/O 区，合起来正好复现上面这张地图。

\begin{pdfdiagram}
  % 地址源与 138 本体（引脚级）
  \node[hw state, minimum width=2.2cm] (addr) at (0.65,1.2) {地址 A[15:0]};
  \node[hw compute, minimum width=2.2cm, minimum height=4.9cm] (hc) at (3.7,-0.2) {74HC138};
  % A[15:13] 接 C/B/A，A12 走底部干线去二级译码
  \draw[hw bus] (addr.east) -- (2.6,1.2) node[pos=0.65,above=1pt,hw label,fill=white,inner sep=1pt]{A14};
  \draw[hw bus] (2.05,1.2) -- (2.05,1.7) -- (2.6,1.7) node[pos=0.6,above=1pt,hw label,fill=white,inner sep=1pt]{A15};
  \draw[hw bus] (2.05,1.2) -- (2.05,0.7) -- (2.6,0.7) node[pos=0.6,above=1pt,hw label,fill=white,inner sep=1pt]{A13};
  \draw[hw bus] (2.05,1.2) -- (2.05,-3.75) -- (6.65,-3.75) -- (6.65,-3.1) node[pos=0.8,above=1pt,hw label,fill=white,inner sep=1pt]{A12};
  \fill[BookMuted] (2.05,1.2) circle (0.05);
  \fill[BookMuted] (2.05,0.7) circle (0.05);
  \node[hw label, anchor=west] at (2.68,1.7) {C};
  \node[hw label, anchor=west] at (2.68,1.2) {B};
  \node[hw label, anchor=west] at (2.68,0.7) {A};
  % 使能引脚：G1=1，/G2A=/G2B=0
  \draw[hw return] (3.0,-3.05) -- (3.0,-2.65);
  \draw[hw return] (3.7,-3.05) -- (3.7,-2.65);
  \draw[hw return] (4.4,-3.05) -- (4.4,-2.65);
  \node[hw label, anchor=south] at (3.0,-2.62) {G1};
  \node[hw label, anchor=south] at (3.7,-2.62) {/G2A};
  \node[hw label, anchor=south] at (4.4,-2.62) {/G2B};
  \node[hw label, anchor=north] at (3.7,-3.12) {G1=1，/G2A=/G2B=0};
  % 输出引脚 Y0--Y7
  \node[hw label, anchor=east] at (4.72,1.925) {Y0};
  \node[hw label, anchor=east] at (4.72,1.375) {Y1};
  \node[hw label, anchor=east] at (4.72,0.825) {Y2};
  \node[hw label, anchor=east] at (4.72,0.275) {Y3};
  \node[hw label, anchor=east] at (4.72,-0.275) {Y4};
  \node[hw label, anchor=east] at (4.72,-0.825) {Y5};
  \node[hw label, anchor=east] at (4.72,-1.375) {Y6};
  \node[hw label, anchor=east] at (4.72,-1.925) {Y7};
  % 片选合并与二级译码
  \node[hw compute, minimum width=1.1cm, minimum height=1.9cm] (or0) at (6.65,1.1) {与门};
  \node[hw compute, minimum width=1.1cm, minimum height=1.65cm] (or1) at (6.65,-0.825) {与门};
  \node[hw state, minimum width=1.5cm] (romcs) at (8.85,1.1) {ROM\_CS};
  \node[hw state, minimum width=1.5cm] (ramcs) at (8.85,-0.825) {RAM\_CS};
  \node[hw compute, minimum width=1.9cm, minimum height=1.1cm] (iodec) at (6.65,-2.55) {二级译码};
  \node[hw state, minimum width=1.7cm] (tail) at (9.45,-2.3) {SRAM 尾窗};
  \node[hw state, minimum width=1.7cm] (iocs) at (9.45,-2.8) {I/O 区 CS};
  \draw[hw bus] (4.8,1.925) -- (6.1,1.925);
  \draw[hw bus] (4.8,1.375) -- (6.1,1.375);
  \draw[hw bus] (4.8,0.825) -- (6.1,0.825);
  \draw[hw bus] (4.8,0.275) -- (6.1,0.275);
  \draw[hw bus] (4.8,-0.275) -- (6.1,-0.275);
  \draw[hw bus] (4.8,-0.825) -- (6.1,-0.825);
  \draw[hw bus] (4.8,-1.375) -- (6.1,-1.375);
  \draw[hw bus] (4.8,-1.925) -- (5.3,-1.925) -- (5.3,-2.55) -- (5.7,-2.55);
  \draw[hw bus] (or0) -- (romcs);
  \draw[hw bus] (or1) -- (ramcs);
  \draw[hw bus] (7.6,-2.3) -- (8.6,-2.3) node[pos=0.5,above=1pt,hw label,fill=white,inner sep=1pt]{A12=0};
  \draw[hw bus] (7.6,-2.8) -- (8.6,-2.8) node[pos=0.5,above=1pt,hw label,fill=white,inner sep=1pt]{A12=1};
\end{pdfdiagram}
\diagramnote{引脚级地址译码。A[15:13] 经 74HC138 译成 8 个 8 KiB 窗口（输出低有效），与门把连续窗口合成 ROM\_CS 与 RAM\_CS，Y7 窗口由 A12 二级译码再分；使能端决定译码器何时工作。}

\begin{longtable}{L{0.09\linewidth}L{0.18\linewidth}L{0.19\linewidth}L{0.44\linewidth}}
\toprule
输出 & C B A（A[15:13]） & 地址区间 & 去向 \\
\midrule
Y0 & 000 & \code{0x0000-0x1FFF} & boot ROM 第 0 页 \\\tablerowrule
Y1 & 001 & \code{0x2000-0x3FFF} & boot ROM 第 1 页 \\\tablerowrule
Y2 & 010 & \code{0x4000-0x5FFF} & boot ROM 第 2 页 \\\tablerowrule
Y3 & 011 & \code{0x6000-0x7FFF} & boot ROM 第 3 页 \\\tablerowrule
Y4 & 100 & \code{0x8000-0x9FFF} & SRAM 第 0 页 \\\tablerowrule
Y5 & 101 & \code{0xA000-0xBFFF} & SRAM 第 1 页 \\\tablerowrule
Y6 & 110 & \code{0xC000-0xDFFF} & SRAM 第 2 页 \\\tablerowrule
Y7 & 111 & \code{0xE000-0xFFFF} & 二级译码：A12=0 归 SRAM 尾窗（\code{0xE000-0xEFFF}），A12=1 为 I/O 区（\code{0xF000-0xFFFF}） \\
\bottomrule
\end{longtable}

这套译码仍然是部分译码：I/O 区里只有 A12 和少数低位参与选择，A[11:5] 等地址位没有进译码器。结果是同一个设备寄存器会在整个 4 KiB 窗口里反复出现——\code{0xF010} 的定时器寄存器，在 \code{0xF110}、\code{0xF210} 等只改变未译码位的地址上同样可见，这些重复地址称为译码别名。SRAM 一侧同理：若实际芯片容量小于所属窗口，未译码的高位也会让同一个存储单元出现在多个地址上。部分译码用更少的门和更短的延迟换来足够大的窗口，在小系统里是合理的工程选择；但它把一道边界推给了规格：别名是否允许必须写明。允许别名，软件就不得用未译码位区分设备，调试时也不能从“访问了哪个地址”反推“选中了哪个设备”；不允许别名，就必须补足低位译码，让未映射组合返回错误或固定值，而不是把一次写错的访问悄悄变成对某个真实设备的读写。

\section{二、valid/ready 握手与背压}

同步互连最小事务可以用 valid/ready 描述：发送者在 valid 为 1 时给出稳定负载，接收者在 ready 为 1 时表示本拍能够接收；只有二者在同一时钟边沿都为 1，事务才真正转移。若 ready 拉低，发送者必须保持 valid 和负载不变，不能悄悄换成下一笔。

\begin{longtable}{L{0.18\linewidth}L{0.28\linewidth}L{0.44\linewidth}}
\toprule
valid & ready & 边沿发生的事 \\
\midrule
0 & 0/1 & 没有有效事务 \\\tablerowrule
1 & 0 & 背压；发送者保持全部负载 \\\tablerowrule
0 & 1 & 接收者空闲等待 \\\tablerowrule
1 & 1 & 当前事务完成一次转移 \\
\bottomrule
\end{longtable}

组合地把 ready 反接到 valid，容易在多个模块串联后形成长组合路径甚至组合环。高频设计会插入 skid buffer 或寄存器切片，在不丢事务的前提下切断这条路径。

\section{三、AXI/APB 的通道与排序边界}

APB 面向低带宽寄存器设备，一次传输用 setup 和 access 阶段完成；AXI 把读地址、读数据、写地址、写数据、写响应分成独立通道，允许多笔事务在途。独立通道意味着写地址与写数据可以不同拍到达，从设备必须分别接收并在配对后执行，不能假设两者同时出现。

事务 ID 用来区分在途请求。互连可以让不同 ID 的响应交错返回，但同一 ID 或带顺序要求的事务必须遵守协议定义。超时、decode error 和 slave error 也要沿响应通道返回；未映射地址不能让总线永远等待。

\begin{keyidea}
地址译码回答“交给谁”，握手回答“何时接收”，ID 与顺序规则回答“哪个响应属于哪笔请求”，错误响应回答“失败怎样闭合”。四者缺一都不是完整互连。
\end{keyidea}

\section{四、仲裁、突发与多笔在途事务}

当 CPU、DMA、显示控制器和网卡都能发起访问时，互连必须先决定谁取得下游端口。仲裁器的输入是一组请求位与当前优先级状态，输出是一组 one-hot 授权位。授权只表示“允许尝试发送”，真正转移仍以 valid/ready 同时有效的采样沿为准。若授权和接收被混为一件事，接收者施加背压时，仲裁器可能提前把授权转给下一方，导致负载与发起者身份错配。

\begin{longtable}{L{0.18\linewidth}L{0.28\linewidth}L{0.32\linewidth}L{0.14\linewidth}}
\toprule
策略 & 状态与选择规则 & 优点与风险 & 适用位置 \\
\midrule
固定优先级 & 始终选择编号最高的有效请求 & 组合路径短；低优先级可能长期无服务 & 紧急控制、实时告警 \\\tablerowrule
轮转 & 每次完成后把指针移到获胜者之后 & 有界等待；需要保存轮转指针 & 通用 DMA 与共享总线 \\\tablerowrule
加权轮转 & 每个主设备按权重消耗服务配额 & 可分配带宽；权重错误会放大尾延迟 & 显示、网络与 CPU 共享内存 \\\tablerowrule
年龄优先 & 选择等待时间最长的事务 & 抑制饥饿；比较网络和时间戳成本较高 & 高并发片上网络 \\
\bottomrule
\end{longtable}

以三个主设备的轮转仲裁为例，A、B、C 同时持续请求。若初始指针指向 A，前三次成功转移依次授权 A、B、C；若 B 在获得授权时遇到下游 \code{ready=0}，B 的授权和负载必须保持，轮转指针也不能前进。指针只能在事务被接受的采样沿更新。这里的状态推进边界与时序逻辑完全一致：组合网络选择候选者，寄存器在成功握手后记录下一次搜索起点。

单拍握手只能搬一个 beat。连续 Cache line、DMA 缓冲区和 DRAM burst 需要把多个 beat 组织成一次突发事务。请求携带起始地址、长度、每 beat 字节数和 burst 类型；接收者用最后一拍标志结束数据流。若 64 B Cache line 经 128-bit 数据通道传输，需要 $64/16=4$ 个 beat。第 2 拍发生背压时，发送者必须保持第 2 拍数据和 byte-enable，不能跳到第 3 拍；最后一拍只有在第 4 个 beat 真正转移时才有效。

多笔在途事务进一步把“提交请求”与“等待响应”分开。发起者为每笔请求分配 ID 或内部表项，保存目标地址、长度、返回缓冲区和错误状态。下游可以并行访问不同 bank，也可以让不同 ID 的响应交错返回。响应到达时，ID 先索引状态表，再把数据写入正确缓冲区；所有 beat 到齐后才能释放表项。状态表已满时，前端必须施加背压，而不是覆盖仍在使用的 ID。

\begin{pdfdiagram}[x=1cm,y=1cm]
  \node[hw state, minimum width=2.4cm] (req) at (0,1.7) {请求入口\\valid/ready};
  \node[hw control, minimum width=2.4cm] (arb) at (3.2,1.7) {仲裁与路由\\选择目标端口};
  \node[hw memory, minimum width=2.4cm] (tab) at (6.4,1.7) {在途表\\ID/长度/状态};
  \node[hw state, minimum width=2.4cm] (dst) at (9.6,1.7) {目标端口\\SRAM/DRAM/设备};
  \node[hw state, minimum width=2.4cm] (rsp) at (6.4,-0.4) {响应入口\\数据/ID/错误};
  \node[hw control, minimum width=2.4cm] (ret) at (3.2,-0.4) {返回路由\\匹配发起者};
  \draw[hw bus] (req) -- (arb) -- (tab) -- (dst);
  \draw[hw return] (dst.south) |- (rsp.east);
  \draw[hw return] (rsp) -- (ret);
  \draw[hw return] (ret.west) -| (req.south);
  \node[hw label] at (4.8,2.25) {接受后分配状态};
  \node[hw label] at (4.8,-0.95) {完成后释放状态};
\end{pdfdiagram}
\diagramnote{多笔在途事务的状态闭环。请求在握手沿被接受后占用一项在途状态；响应按 ID 返回并累计 beat，只有数据完整或错误被确认后才释放。请求、接受、响应与完成是四个不同边界。}

\section{五、协议桥、时钟域与宽度转换}

片上系统不会让所有设备使用同一协议、位宽和时钟。高性能主干可能使用 256-bit AXI，低速寄存器设备使用 32-bit APB，传感器控制器还可能运行在更慢的独立时钟域。协议桥的职责不是改接线名称，而是保存尚未完成的事务，并在两侧分别遵守各自契约。

AXI 到 APB 的写事务需要先分别接收写地址和写数据，确认两者属于同一事务，再在 APB 侧产生 setup 与 access 两阶段。APB 从设备拉长 \code{PREADY} 时，桥必须保持地址、写数据和控制位；APB 完成后，桥才在 AXI 写响应通道返回状态。AXI 的写请求被桥接收并不表示 APB 寄存器已经更新，写响应才是跨过低速设备边界的完成点。

位宽转换也包含状态。128-bit 主干向 32-bit 外设窗口写 16 B 时，桥要拆成四笔 32-bit 访问，并按 byte strobe 处理首尾部分写；反方向读取时要把四个返回字组合后再交付上游。任何子访问失败，桥都必须结束剩余动作并返回可定位的错误，不能用未读取的数据填满结果后报告成功。

跨时钟域不能直接把多位 payload 的每一位各自接两级同步器。单 bit 状态可以同步，但多位数据的各 bit 可能在不同目标时钟沿被采到。常用方法是异步 FIFO：写指针在写时钟域推进，读指针在读时钟域推进；指针转成 Gray code 后跨域同步，因为相邻 Gray 值只改变一位。满和空由同步后的对侧指针计算，数据本体写入双口存储阵列。复位时两侧必须分别回到已知指针，并在同步延迟结束前禁止误报可读数据。

\begin{longtable}{L{0.2\linewidth}L{0.30\linewidth}L{0.30\linewidth}L{0.12\linewidth}}
\toprule
转换 & 必须保存的状态 & 完成边界 & 典型错误 \\
\midrule
AXI 到 APB & 地址、写数据、byte strobe、响应状态 & APB access 被接受并形成 AXI 响应 & 把 AXI 接受误当作设备完成 \\\tablerowrule
宽到窄 & 子访问序号、累计数据、首个错误 & 全部子访问完成或首个错误闭合 & 部分写覆盖邻接字节 \\\tablerowrule
窄到宽 & 聚合缓冲、有效字节掩码 & 满足上游 beat 宽度与有效掩码 & 未初始化字节被当成有效数据 \\\tablerowrule
跨时钟域 & Gray 指针、同步级、满空状态 & 目标域在本地时钟沿确认可见 & 多位总线逐位同步造成混合值 \\
\bottomrule
\end{longtable}

\section{六、Crossbar、NoC 与一致性互连}

单共享总线在任一时刻只能服务一组主从传输。主设备和从设备增加后，crossbar 允许访问不同目标的事务并行通过：CPU 访问 SRAM 的同时，DMA 可以访问 DRAM；只有目标相同时才在对应输出端仲裁。代价是路由、仲裁和返回路径随端口数增加，面积和布线迅速增长。

更大的 SoC 使用片上网络（Network-on-Chip, NoC）。请求被封装成若干 flit，经路由器逐跳前进。每个输入端口按虚通道保存队列，路由计算选择下一跳，交换分配决定本拍哪个 flit 穿过交换结构。credit 表示下游还剩多少缓冲槽；发送一个 flit 消耗一个 credit，下游释放槽位后返还 credit。credit 为零时必须停止发送，否则下游缓冲会溢出。

NoC 的错误不仅是数据损坏，还包括协议死锁。若请求通道等待响应通道释放，而响应通道又等待请求通道持有的资源，所有节点都可能保持有效但永远不前进。工程上用独立虚通道、规定的通道依赖顺序和超时检测打破环形等待。QoS 只决定谁先得到带宽，不能修复通道依赖形成的死锁。

一致性互连还要维护 Cache line 的共享状态。一次核心写 miss 可能发出 Read-for-Ownership，请求目录定位其他副本并向持有者发送失效消息；所有失效确认返回后，写者才能取得独占所有权。目录项、探测请求、确认和数据返回分别占用不同通道。若互连只交付数据却遗漏失效确认，两个核心可能同时保留可写副本，错误会在软件层表现为偶发的数据竞争。

\section{七、一次事务怎样成功或失败地闭合}

考虑 CPU 对设备窗口的一次 32-bit MMIO 读。地址先经过译码和权限检查，crossbar 为目标端口仲裁，协议桥把请求转换为设备侧访问，设备状态机在若干周期后给出数据或错误，返回网络按事务 ID 把响应送回 CPU。CPU 只有收到响应后才能让该 load 在体系结构上完成。中间任一环节都必须给出有限结局：接受、背压、错误或超时。

\begin{longtable}{L{0.18\linewidth}L{0.30\linewidth}L{0.30\linewidth}L{0.14\linewidth}}
\toprule
故障位置 & 可观察症状 & 必须保留的诊断状态 & 恢复动作 \\
\midrule
地址译码 & 立即返回 decode error & 请求地址、发起者 ID & 修正地址地图或驱动 \\\tablerowrule
从设备 & slave error 或设备状态错误 & 设备寄存器、事务类型 & 清错、复位设备或终止请求 \\\tablerowrule
链路背压 & 延迟增加但事务仍推进 & 队列水位、ready 停顿周期 & 降低注入率或增加缓冲 \\\tablerowrule
响应丢失 & 在途表项长时间不释放 & ID、目标端口、超时计数 & 隔离端口并按协议复位 \\\tablerowrule
一致性探测 & 所有权请求等待确认 & 目录状态、待确认位图 & 处理失联核心或机器检查 \\
\bottomrule
\end{longtable}

\section*{章末练习}

\begin{chapterexercise}{握手概念}
解释请求有效、被接受、返回响应和事务完成四个边界
\exercisecheck{检查要点}{不把 valid 单独有效误认为已经转移}
\end{chapterexercise}

\begin{chapterexercise}{轮转 trace}
三个主设备持续请求，写出六次成功握手的授权顺序，并加入一次两拍背压
\exercisecheck{检查要点}{指针只在成功握手后推进}
\end{chapterexercise}

\begin{chapterexercise}{突发边界}
64 B Cache line 经 128-bit 通道传输，计算 beat 数并说明第 2 beat 背压时各信号行为
\exercisecheck{检查要点}{数据、strobe、last 都必须保持对应关系}
\end{chapterexercise}

\begin{chapterexercise}{协议桥}
画出一笔 AXI 写经 APB 桥到设备寄存器的时间线
\exercisecheck{检查要点}{地址与数据可分开到达，响应只能在 APB 完成后返回}
\end{chapterexercise}

\begin{chapterexercise}{跨域排错}
一个 8-bit 计数值逐位用同步器跨域，说明为何读到不存在的中间值并给出替代方案
\exercisecheck{检查要点}{多位数据使用握手或异步 FIFO}
\end{chapterexercise}

\begin{chapterexercise}{故障定位}
MMIO 读超时，按译码、仲裁、桥、从设备和返回路由列出测点
\exercisecheck{检查要点}{每层都要有前进状态或有限错误}
\end{chapterexercise}

\section*{参考解答与要点}

\begin{chapteranswer}{握手概念}
valid 表示负载有效，ready 表示接收能力；二者在同一采样沿有效才算接受。响应可能更晚到达，完整数据或错误返回后事务才完成。
\answercheck{必须保留的检查点}{请求、接受、响应、完成分开表述}
\end{chapteranswer}

\begin{chapteranswer}{轮转 trace}
无背压时顺序为 A、B、C、A、B、C。若 B 获授权后 ready 连续两拍为 0，B 保持授权和负载；第 3 拍成功后才转向 C。
\answercheck{必须保留的检查点}{背压期间不能改变发起者或 payload}
\end{chapteranswer}

\begin{chapteranswer}{突发边界}
每 beat 16 B，共 4 beat。第 2 beat 背压时发送者保持第 2 beat 数据、strobe 与 valid；接受后才推进到第 3 beat，第 4 beat 才置 last。
\answercheck{必须保留的检查点}{last 属于具体 beat，不能按墙钟拍数提前}
\end{chapteranswer}

\begin{chapteranswer}{协议桥}
AXI 分别接收 AW 与 W；桥配对后发 APB setup，下一拍进入 access，并在 PREADY 有效沿采样结果；随后产生 B 响应。
\answercheck{必须保留的检查点}{AXI 入口接受不等于设备寄存器已更新}
\end{chapteranswer}

\begin{chapteranswer}{跨域排错}
各 bit 的同步延迟可能不同，目标域会组合出源域从未产生的值。稳定快照可用 request/ack 握手；连续流使用带 Gray 指针的异步 FIFO。
\answercheck{必须保留的检查点}{两级同步器只直接适用于单 bit 控制}
\end{chapteranswer}

\begin{chapteranswer}{故障定位}
先查地址是否命中唯一目标，再查请求是否获授权、桥两侧握手是否推进、设备状态机是否结束、响应 ID 是否返回。任一点停止都应有计数或超时记录。
\answercheck{必须保留的检查点}{不从 CPU 超时直接推断设备损坏}
\end{chapteranswer}
